\documentclass{amsart}
% For dealing with references we use the comment environment
\usepackage{verbatim}
\newenvironment{reference}{\comment}{\endcomment}
%\newenvironment{reference}{}{}
\newenvironment{slogan}{\comment}{\endcomment}
\newenvironment{history}{\comment}{\endcomment}

% For commutative diagrams we use Xy-pic
\usepackage[all]{xy}

% We use 2cell for 2-commutative diagrams.
\xyoption{2cell}
\UseAllTwocells

% We use multicol for the list of chapters between chapters
\usepackage{multicol}

% This is generally recommended for better output
\usepackage{lmodern}
\usepackage[T1]{fontenc}

% For cross-file-references

% Package for hypertext links:
\usepackage{hyperref}

% For any local file, say "hello.tex" you want to link to please

% Theorem environments.
%
\theoremstyle{plain}
\newtheorem{theorem}[subsection]{Theorem}
\newtheorem{proposition}[subsection]{Proposition}
\newtheorem{lemma}[subsection]{Lemma}

\theoremstyle{definition}
\newtheorem{definition}[subsection]{Definition}
\newtheorem{example}[subsection]{Example}
\newtheorem{exercise}[subsection]{Exercise}
\newtheorem{situation}[subsection]{Situation}

\theoremstyle{remark}
\newtheorem{remark}[subsection]{Remark}
\newtheorem{remarks}[subsection]{Remarks}

\numberwithin{equation}{subsection}

% Macros
%
\def\lim{\mathop{\mathrm{lim}}\nolimits}
\def\colim{\mathop{\mathrm{colim}}\nolimits}
\def\Spec{\mathop{\mathrm{Spec}}}
\def\Hom{\mathop{\mathrm{Hom}}\nolimits}
\def\Ext{\mathop{\mathrm{Ext}}\nolimits}
\def\SheafHom{\mathop{\mathcal{H}\!\mathit{om}}\nolimits}
\def\SheafExt{\mathop{\mathcal{E}\!\mathit{xt}}\nolimits}
\def\Sch{\mathit{Sch}}
\def\Mor{\mathop{\mathrm{Mor}}\nolimits}
\def\Ob{\mathop{\mathrm{Ob}}\nolimits}
\def\Sh{\mathop{\mathit{Sh}}\nolimits}
\def\NL{\mathop{N\!L}\nolimits}
\def\CH{\mathop{\mathrm{CH}}\nolimits}
\def\proetale{{pro\text{-}\acute{e}tale}}
\def\etale{{\acute{e}tale}}
\def\QCoh{\mathit{QCoh}}
\def\Ker{\mathop{\mathrm{Ker}}}
\def\Im{\mathop{\mathrm{Im}}}
\def\Coker{\mathop{\mathrm{Coker}}}
\def\Coim{\mathop{\mathrm{Coim}}}

% Boxtimes
%
\DeclareMathSymbol{\boxtimes}{\mathbin}{AMSa}{"02}

%
% Macros for moduli stacks/spaces
%
\def\QCohstack{\mathcal{QC}\!\mathit{oh}}
\def\Cohstack{\mathcal{C}\!\mathit{oh}}
\def\Spacesstack{\mathcal{S}\!\mathit{paces}}
\def\Quotfunctor{\mathrm{Quot}}
\def\Hilbfunctor{\mathrm{Hilb}}
\def\Curvesstack{\mathcal{C}\!\mathit{urves}}
\def\Polarizedstack{\mathcal{P}\!\mathit{olarized}}
\def\Complexesstack{\mathcal{C}\!\mathit{omplexes}}
% \Pic is the operator that assigns to X its picard group, usage \Pic(X)
% \Picardstack_{X/B} denotes the Picard stack of X over B
% \Picardfunctor_{X/B} denotes the Picard functor of X over B
\def\Pic{\mathop{\mathrm{Pic}}\nolimits}
\def\Picardstack{\mathcal{P}\!\mathit{ic}}
\def\Picardfunctor{\mathrm{Pic}}
\def\Deformationcategory{\mathcal{D}\!\mathit{ef}}

\setlength{\emergencystretch}{3em}
\begin{document}
\title[Stacks Project, Tag 0GU8]{446 --- Chow homology descent exact sequence for proper completely decomposed envelopes}
\maketitle
\noindent Copyright (C) 2005--2025 Johan de Jong. Stacks Project authors. Source: \href{https://stacks.math.columbia.edu/tag/0GU8}{Tag 0GU8}.

\noindent Editorial difficulty note (not author text). Difficulty H2: Surjectivity requires locally finite birational lifts of all cycle supports, including supports of rational-equivalence relations. The kernel must then be represented by cycles on the double fibre product, with explicit generic components and degree cancellation distinguishing equal-dimensional and lower-dimensional images..

\noindent Editorial provenance note (not author text). History: excerpt from revision \texttt{a0444\allowbreak 6e57e\allowbreak c1fbc\allowbreak 252a8\allowbreak 71afc\allowbreak ec775\allowbreak 2fb28\allowbreak 07b14}, chow.tex, lines 3592--3709. Reference presentation was adapted; original-author context and core support blocks were retained or added. The mathematical wording is unchanged. Licensed under GNU FDL 1.2 or later; no invariant sections or cover texts. See ../licenses/COPYING. Context blocks reproduce author definitions and supporting results; they are not additional counted problems.

\begin{definition}
\label{definition-envelope}
\begin{reference}
\href{https://stacks.math.columbia.edu/bibliography}{Bibliography: F (Definition 18.3)}
\end{reference}
Let $X$ be a scheme. An {\it envelope} is a proper morphism $f : Y \to X$
which is completely decomposed
(More on Morphisms, Definition \ref{more-morphisms-definition-cd-morphism}).
\end{definition}

\begin{situation}
\label{situation-setup}
Here $S$ is a locally Noetherian, and universally catenary scheme.
Moreover, we assume $S$ is endowed with a dimension function
$\delta : S \longrightarrow \mathbf{Z}$.
\end{situation}

\begin{definition}
\label{more-morphisms-definition-cd-morphism}
A morphism $f : X \to Y$ of schemes is said to be
{\it completely decomposed}\footnote{This may be nonstandard terminology.}
if for all points $y \in Y$ there
is a point $x \in X$ with $f(x) = y$ such that the field
extension $\kappa(x)/\kappa(y)$ is trivial.
A family of morphisms $\{f_i : X_i \to Y\}_{i \in I}$ of
schemes with fixed target is said to be {\it completely decomposed}
if $\coprod f_i : \coprod Y_i \to X$ is completely decomposed.
\end{definition}

\begin{lemma}
\label{lemma-envelope}
Let $(S, \delta)$ be as in Situation \ref{situation-setup}.
Let $X$ be a scheme locally of finite type over $S$.
Let $f : Y \to X$ be an envelope. Then
we have an exact sequence
$$
\CH_k(Y \times_X Y) \xrightarrow{p_* - q_*}
\CH_k(Y) \xrightarrow{f_*}
\CH_k(X) \to 0
$$
for all $k \in \mathbf{Z}$. Here $p, q : Y \times_X Y \to Y$ are
the projections.
\end{lemma}

\begin{proof}
Since $f$ is an envelope, $f$ is proper and hence pushforward on
cycles and cycle classes is defined, see
Sections \href{https://stacks.math.columbia.edu/tag/02R3}{section proper pushforward (Tag 02R3)} and \href{https://stacks.math.columbia.edu/tag/02RF}{section push pull (Tag 02RF)}.
Similarly, the morphisms $p$ and $q$ are proper as base changes of $f$.
The composition of the arrows is zero as
$f_* \circ p_* = (p \circ f)_* = (q \circ f)_* = f_* \circ q_*$, see
Lemma \href{https://stacks.math.columbia.edu/tag/02R5}{lemma compose pushforward (Tag 02R5)}.

\medskip\noindent
Let us show that $f_* : Z_k(Y) \to Z_k(X)$ is surjective.
Namely, suppose that we have $\alpha = \sum n_i[Z_i] \in Z_k(X)$
where $Z_i \subset X$ is a locally finite family of integral
closed subschemes. Let $x_i \in Z_i$ be the generic point.
Since $f$ is an envelope and hence completely decomposed,
there exists a point $y_i \in Y$ with $f(y_i) = x_i$
and with $\kappa(y_i)/\kappa(x_i)$ trivial. Let $W_i \subset Y$
be the integral closed subscheme with generic point $y_i$.
Since $f$ is closed, we see that $f(W_i) = Z_i$.
It follows that the family of closed subschemes $W_i$ is locally finite
on $Y$. Since $\kappa(y_i)/\kappa(x_i)$ is trivial we see
that $\dim_\delta(W_i) = \dim_\delta(Z_i) = k$. Hence
$\beta = \sum n_i[W_i]$ is in $Z_k(Y)$. Finally, since
$\kappa(y_i)/\kappa(x_i)$ is trivial, the degree of the dominant
morphism $f|_{W_i} : W_i \to Z_i$ is $1$ and we conclude
that $f_*\beta = \alpha$.

\medskip\noindent
Since $f_* : Z_k(Y) \to Z_k(X)$ is surjective, a fortiori the map
$f_* : \CH_k(Y) \to \CH_k(X)$ is surjective.

\medskip\noindent
Let $\beta \in Z_k(Y)$ be an element such that $f_*\beta$ is zero in
$\CH_k(X)$. This means we can find a locally finite family of
integral closed subschemes $Z_j \subset X$ with $\dim_\delta(Z_j) = k + 1$
and $f_j \in R(Z_j)^*$ such that
$$
f_*\beta = \sum (Z_j \to X)_*\text{div}(f_j)
$$
as cycles where $i_j : Z_j \to X$ is the given closed immersion.
Arguing exactly as above, we can find a locally finite
family of integral closed subschemes $W_j \subset Y$
with $f(W_j) = Z_j$ and such that $W_j \to Z_j$ is birational, i.e.,
induces an isomorphism $R(Z_j) = R(W_j)$. Denote $g_j \in R(W_j)^*$
the element corresponding to $f_j$. Observe that $W_j \to Z_j$
is proper and that $(W_j \to Z_j)_*\text{div}(g_j) = \text{div}(f_j)$
as cycles on $Z_j$. It follows from this that if we replace
$\beta$ by the rationally equivalent cycle
$$
\beta' = \beta - \sum (W_j \to Y)_*\text{div}(g_j)
$$
then we find that $f_*\beta' = 0$.
(This uses Lemma \href{https://stacks.math.columbia.edu/tag/02R5}{lemma compose pushforward (Tag 02R5)}.)
Thus to finish the proof
of the lemma it suffices to show the claim in the following paragraph.

\medskip\noindent
Claim: if $\beta \in Z_k(Y)$ and $f_*\beta = 0$, then
$\beta = \delta + p_*\gamma - q_*\gamma$ in $Z_k(Y)$ for some
$\gamma \in Z_k(Y \times_X Y)$. Namely, write $\beta = \sum_{j \in J} n_j[W_j]$
with $\{W_j\}_{j \in J}$ a locally finite family of integral closed
subschemes of $Y$ with $\dim_\delta(W_j) = k$.
Fix an integral closed subscheme $Z \subset X$. Consider the subset
$J_Z = \{j \in J : f(W_j) = Z\}$. This is a finite set. There are three
cases:
\begin{enumerate}
\item $J_Z = \emptyset$. In this case we set $\gamma_Z = 0$.
\item $J_Z \not = \emptyset$ and $\dim_\delta(Z) = k$.
The condition $f_*\beta = 0$ implies by looking at the
coefficient of $Z$ that $\sum_{j \in J_Z} n_j\deg(W_j/Z) = 0$.
In this case we choose an integral closed subscheme $W \subset Y$
which maps birationally onto $Z$ (see above). Looking at generic
points, we see that $W_j \times_Z W$ has a unique irreducible
component $W'_j \subset W_j \times_Z W \subset Y \times_X Y$
mapping birationally to $W_j$. Then $W'_j \to W$ is dominant
and $\deg(W'_j/W) = \deg(W_j/W)$. Thus if we set
$\gamma_Z = \sum_{j \in J_Z} n_j[W'_j]$
then we see that
$p_*\gamma_Z = \sum_{j \in J_Z} n_j[W_j]$ and
$q_*\gamma_Z = \sum_{j \in J_Z} n_j\deg(W'_j/W)[W] = 0$.
\item $J_Z \not = \emptyset$ and $\dim_\delta(Z) < k$.
In this case we choose an integral closed subscheme $W \subset Y$
which maps birationally onto $Z$ (see above). Looking at generic
points, we see that $W_j \times_Z W$ has a unique irreducible
component $W'_j \subset W_j \times_Z W \subset Y \times_X Y$
mapping birationally to $W_j$. Then $W'_j \to W$ is dominant
and $k = \dim_\delta(W'_j) > \dim_\delta(W) = \dim_\delta(Z)$.
Thus if we set $\gamma_Z = \sum_{j \in J_Z} n_j[W'_j]$
then we see that
$p_*\gamma_Z = \sum_{j \in J_Z} n_j[W_j]$ and
$q_*\gamma_Z = 0$.
\end{enumerate}
Since the family of integral closed subschemes $\{f(W_j)\}$
is locally finite on $X$
(Lemma \href{https://stacks.math.columbia.edu/tag/02R2}{lemma quasi compact locally finite (Tag 02R2)})
we see that the $k$-cycle
$$
\gamma = \sum\nolimits_{Z \subset X\text{ integral closed}} \gamma_Z
$$
on $Y \times_X Y$ is well defined. By our computations above it follows that
$p_*\gamma_Z = \beta$ and $q_*\gamma_Z = 0$ which implies
what we wanted to prove.
\end{proof}
\end{document}
